%=================================================================
% Cover Letter for submission to MDPI Big Data and Cognitive Computing (BDCC)
% Special Issue: "Natural Language Processing Applications in Big Data"
% Manuscript: UzABSA-LLM: Parameter-Efficient Fine-Tuning and
%             Multi-Domain Annotation for Uzbek Aspect-Based Sentiment Analysis
%=================================================================
\documentclass[11pt]{letter}

\usepackage[a4paper,margin=1in]{geometry}
\usepackage[T1]{fontenc}
\usepackage[utf8]{inputenc}
\usepackage{microtype}
\usepackage{hyperref}
\hypersetup{hidelinks}

\setlength{\parindent}{0pt}
\setlength{\parskip}{0.6em}

\signature{%
  Sanatbek Matlatipov\\
  on behalf of all co-authors\\
  National University of Uzbekistan named after Mirzo Ulugbek\\
  Tashkent 100174, Uzbekistan\\
  Email: s.matlatipov@nuu.uz}

\address{%
  Sanatbek Matlatipov (Corresponding Author)\\
  Faculty of Applied Mathematics and Intelligent Technologies\\
  National University of Uzbekistan named after Mirzo Ulugbek\\
  Tashkent 100174, Uzbekistan}

\begin{document}

\begin{letter}{%
  The Editorial Office\\
  \textit{Big Data and Cognitive Computing} (MDPI)\\
  Special Issue: ``Natural Language Processing Applications in Big Data''}

\opening{Dear Editors,}

We submit our manuscript \textbf{``UzABSA-LLM: Parameter-Efficient Fine-Tuning and Multi-Domain Annotation for Uzbek Aspect-Based Sentiment Analysis''} for consideration as an original research article in \textit{Big Data and Cognitive Computing}, for the Special Issue ``Natural Language Processing Applications in Big Data''.

\textbf{Significance.}
Aspect-Based Sentiment Analysis (ABSA) is well established for high-resource languages, but Uzbek---a Turkic language spoken by over 38 million people---remains under-served owing to its sparse representation in multilingual pre-training corpora, its agglutinative morphology, and its Latin/Cyrillic script duality. To the best of our knowledge, this is the first controlled comparison of QLoRA-adapted open-weight 7--8B language models on the UzABSA benchmark, benchmarked against monolingual Uzbek encoders. We reformulate ABSA as a single-pass generative task with Uzbek instruction tuning, so that a model extracts aspect terms and assigns their polarity in one structured JSON generation that feeds a downstream annotation pipeline.

\textbf{Findings in context.}
Generative ABSA has concentrated on English, and prior Uzbek work has relied on rule-based methods, classical machine learning, or encoder models. Extending the generative paradigm to a low-resource, morphologically rich language, we report: (i)~a controlled comparison in which Qwen~2.5-7B attains the best extraction F1 (0.708) at a 100\% parse rate and Llama~3.1-8B the best end-to-end pair F1 (0.651) and sentiment accuracy (0.925), yet a paired bootstrap detects no significant difference between them---reported as a non-result rather than a declared winner, under a multiple-comparison correction applied to the full family of tests; (ii)~a cost--capability trade-off in which fine-tuned 110M-parameter Uzbek encoders show no statistically detectable disadvantage at extraction while the larger models are clearly stronger at sentiment and joint prediction; (iii)~zero-shot controls showing that QLoRA fine-tuning at least doubles exact extraction F1 and triples pair F1, buying task competence and prediction calibration rather than format compliance; and (iv)~a validated model-assisted annotation pipeline applied to 5{,}038 reviews across 23 domains, yielding 9{,}884 aspect-level annotations with per-item quality metadata. All code, data, and models are publicly released.

\textbf{Fit to BDCC.} The work speaks to both halves of the journal's remit.

\emph{Big data.} For user-generated text in a low-resource language, the binding constraint on large-scale analysis is not storage or compute but \emph{annotation}: a heterogeneous review stream cannot become an analytical resource without labels that do not exist and cannot be hand-produced at scale. We contribute a pipeline that performs that conversion, and we report its economics from measured runs rather than asserting scalability---annotation proceeds at roughly 840 reviews per hour on a single GPU, and automated quality scoring of the entire corpus costs under one US dollar. The dominant remaining cost is human, which is why the automatic layers are built to be applied exhaustively while scarce native-speaker effort is spent calibrating them. The pipeline is model-agnostic and transfers across languages and domains; our 23-domain corpus exhibits the heterogeneity typical of real review streams, and we measure how annotation quality degrades with distance from the training domain.

\emph{Cognitive computing.} The pipeline's second layer substitutes machine judgement for human judgement, and we treat that substitution as a claim to be tested rather than assumed. Calibrating an LLM judge against native-speaker ratings, we report both where it succeeds (Spearman $\rho=0.81$ on sentiment, $\rho=0.66$ on the tier-defining overall score) and where it does not (weaker agreement on completeness, where it is systematically harsher than humans). The accompanying validation study yields a result about the task itself: polarity judgement proved essentially unambiguous---annotators agreed on the sentiment of every jointly identified aspect---whereas \emph{what counts as the aspect span} is the dominant source of disagreement even between trained native speakers. Human and machine annotators diverge in the same place and for the same reason, and it is this calibration that licenses the quality tiers on which the released dataset depends.

\textbf{Prior MDPI submission.}
As required, we disclose that an earlier version of this manuscript was submitted to \textit{Computers} (MDPI) under manuscript ID \textbf{computers-4419237} and was not accepted; this ID is also supplied in the submission system. The present version differs substantially: all models were retrained and re-evaluated after we identified and corrected two defects in the earlier pipeline (a system-prompt serialization bug and text-level overlap between the training and evaluation splits), and we added zero-shot controls, encoder baselines, paired-bootstrap significance testing across all system pairs, and a native-speaker validation study. The reported results are new, and the claims have been narrowed to what the corrected experiments support. The author list is unchanged from that submission.

We confirm that neither the manuscript nor any parts of its content are currently under consideration for publication with or published in another journal. All authors have approved the manuscript and agree with its submission to BDCC.

Thank you for considering our manuscript. We would be happy to provide any further information required.

\closing{Sincerely,}

\end{letter}
\end{document}
